% Generated draft: requires mathematical and source review.
\documentclass[11pt]{article}
\usepackage[margin=25mm]{geometry}
\usepackage{fontspec}
\setmainfont{Latin Modern Roman}
\usepackage{amsmath,amssymb,mathtools,unicode-math}
\setmathfont{Latin Modern Math}
\usepackage{enumitem,needspace,xurl,hyperref}
\setlength{\parindent}{0pt}
\setlength{\parskip}{5pt}
\setlength{\emergencystretch}{3em}
\newcommand{\N}{\mathbb{N}}
\newcommand{\Z}{\mathbb{Z}}
\newcommand{\Q}{\mathbb{Q}}
\newcommand{\R}{\mathbb{R}}
\newcommand{\C}{\mathbb{C}}
\newcommand{\F}{\mathbb{F}}
\newcommand{\A}{\mathbb{A}}
\newcommand{\PP}{\mathbb{P}}
\newcommand{\E}{\mathbb{E}}
\newcommand{\eps}{\varepsilon}
\newcommand{\dd}{\,\mathrm{d}}
\begin{document}
\section*{UnsolvedMath ID: 701395}
\textbf{Source problem code:} catalogue entry
\subsection{Definitions and mathematical statement}
Let \(G\) be a compact topological group with identity \(e\). For \(g\in G\), define \(\operatorname{ord}(g)\) to be the least positive integer \(m\) such that \(g^m=e\), if such an integer exists, and set \(\operatorname{ord}(g)=\infty\) otherwise. Regard \(\infty\) as larger than every positive integer.

Assume that the orders of elements of \(G\) are unbounded, meaning
\[
\forall N\in\Z_{\geq 1}\;\exists g\in G\quad \operatorname{ord}(g)>N.
\]
Must there exist \(h\in G\) such that
\[
\operatorname{ord}(h)=\infty,
\qquad\text{equivalently}\qquad
\forall m\in\Z_{\geq 1},\quad h^m\ne e?
\]
The source does not explicitly specify whether Hausdorffness is included in its convention for a topological group; that convention requires confirmation. If ``arbitrarily high orders'' specifically means arbitrarily large finite orders, the resulting universal question is equivalent: any counterexample to either formulation would have only finite-order elements.
\subsection{Short English statement}
Must every compact topological group with unbounded element orders contain an element of infinite order? The supplied record calls this a well-known problem and attributes its proposal to J. Mycielski in 2010.
\subsection{Sources}
\begin{itemize}
\item Local UnsolvedMath registry, supplied source record: permanent ID \(701395\), exact source code \texttt{catalogue entry}, title \emph{Kourovka Notebook Problem 17.93}. The supplied registry labels the problem open; this status has not been independently verified.
\item \emph{The Kourovka Notebook (Unsolved Problems in Group Theory)}, 17th issue (2010), Problem 17.93; proposed by J. Mycielski. The supplied statement describes it as a well-known problem.
\item The supplied background also identifies the 21st issue edition (2026), \texttt{arXiv:1401.0300v48}. The supplied collection description names E. I. Khukhro and V. D. Mazurov as editors. These bibliographic details are supplied attribution and have not been independently checked.
\end{itemize}
\end{document}
